\documentclass[10pt,a4paper]{amsart}
\usepackage[margin=19mm]{geometry}
\usepackage{amsmath,amssymb,mathtools}
\usepackage[scr=rsfs,cal=euler]{mathalfa}
\usepackage{xfrac}
\usepackage[unicode,hidelinks,bookmarks=false]{hyperref}
\newcommand{\ie}{i.e.\ }
\newcommand{\cf}{cf.\ }
\newcommand{\resp}{respectively\ }
\newcommand{\Subsection}{\subsection}
\providecommand{\nobreakdash}{\nobreak\hbox{-}}
\setlength{\emergencystretch}{2em}
\multlinegap0pt
\usepackage[shortlabels]{enumitem}
\usepackage{mdframed}
\usepackage{amssymb}
\newtheorem{theorem}{Theorem}
\newtheorem{lemma}{Lemma}
\newcommand{\norm}[1]{\left\lVert #1 \right\rVert}
\title{Non-uniform Kahn-Kalai, spread, variants, and applications}
\author{Thinula De Silva, Pu Gao}
\date{Source: arXiv:2603.13737v1 (CC BY 4.0)}
\begin{document}
\maketitle

\noindent\emph{Source and licence.} Non-uniform Kahn-Kalai, spread, variants, and applications. Version arXiv:2603.13737v1 (CC BY 4.0), distributed by arXiv under the Creative Commons Attribution 4.0 International licence, http://creativecommons.org/licenses/by/4.0/, as recorded in the arXiv licence metadata for this exact version and shown on the abstract page https://arxiv.org/abs/2603.13737v1. Full licence notice: \texttt{licenses/arxiv-2603.13737-CC-BY-4.0.txt}.\par\smallskip

\noindent\textit{Editorial scope.} Counted unit: the article's lemma Lemma (source0.tex lines 1022--1025). The article's complete original proof of it is delivered with it (source0.tex lines 1027--1061, one \texttt{proof} environment of 35 lines). No mathematics is written, repaired or re-derived: the statement and the proof are the article's own lines, in the article's own notation, and no retained line is substituted or re-flowed.  Retained uncounted as the source-local context the proof needs: lines 378--389.  Nothing was omitted from any retained span, so the package carries no omission record.  No retained line contains a citation, so the wrapper carries no bibliography and no bibliography entry.  No retained line contains a figure environment, an image-inclusion command or drawing code, so no image is delivered and the package declares no asset dependency.  Editorial formatting only, in addition: an AMS \texttt{amsart} preamble that restates the article's own declarations and macros used by the retained lines, namely the environments \texttt{theorem}, \texttt{lemma} on separate counters, together with the standard packages those lines need.  The retained environments print in this document as Theorem 1, Lemma 1 in order of appearance, whereas the article prints the same items with the numbering of its own full document.  The complete native source of the article as distributed in the arXiv e-print for this exact version is bundled unchanged as \texttt{sources/196349/source0.tex}.
\begin{theorem}
\label{theorem: G(n,d) bi-valued 0-statement}
      Let $\mathcal{D}$ be the family of bi-valued $\vec{d}$  with $\norm{\vec{d}}_1=0\pmod{2}$ and the following properties:
      \begin{enumerate}[label=(A\arabic*), leftmargin=3cm]
          \item $d_2^2 = o(n_1)$ \label{cond: d_2 and n_1 bound}
          \item $d_1^2 = o(\sqrt{n d_2})$ \label{cond: stronger d_1 bound}
          \item $d_2\gg\log n$ \label{cond: d_2 lower bound embedding}
          \item $n_2^2 d_2^3 = o(\norm{\vec{d}}_1^2)$ \label{cond: n_2 and d_2 bound with norm}
          \item $n_1\leq n^\delta$ for some fixed $\delta < 1$ \label{cond: n_1 small}
      \end{enumerate}
      Then, both the 1-statement and ideal 0-statement hold for $\mathcal{D}$.
\end{theorem}

\begin{lemma}
\label{lemma: G(n,d) 0-statement Intermediate Lemma}
Let $A = \sum\limits_{i = 1}^{n} d_i(d_i-1)$, $T= O(d_1^2)$, and $B = A - T$. Furthermore, let $A' = A/(2\norm{\vec{d}}_1)$ and $B' = B/(2\norm{\vec{d}}_1 - k d_2)$ for some $k = \Theta(1)$. If $Z = (A' - B')(1 + A' + B')$, then $Z = o(1)$.
\end{lemma}

\begin{proof}
Note that
\[
A' - B' = \frac{A}{2\norm{\vec{d}}_1} - \frac{A-T}{2\norm{\vec{d}}_1 - k d_2} = \frac{2T\norm{\vec{d}}_1 - kA d_2}{2\norm{\vec{d}}_1(2\norm{\vec{d}}_1 - k d_2)} = \frac{T}{2\norm{\vec{d}}_1-k d_2} - A'\cdot\frac{k d_2}{\norm{\vec{d}}_1-k d_2},
\]
and so it follows that
\[
A' + B' = A' + \left(1 +\frac{k d_2}{\norm{\vec{d}}_1-k d_2}\right)A' - \frac{T}{2\norm{\vec{d}}_1-k d_2} = \frac{2\norm{\vec{d}}_1 - k d_2}{\norm{\vec{d}}_1 - k d_2}\cdot A' - \frac{T}{2\norm{\vec{d}}_1-k d_2}.
\]
Note that $A' = \Theta((n_1 d_1^2 + n_2 d_2^2)/(n_1 d_1 + n_2 d_2))= \Omega(1)$. Furthermore, since $n_1 d_1 \leq n^{\delta}\sqrt{d_2} = o(n_2 d_2)$ by conditions~\ref{cond: stronger d_1 bound} and ~\ref{cond: n_1 small}, we have
\begin{equation}
\label{eqn: A' bound}
    A' = O\left(\frac{n_1 d_1^2 +n_2 d_2^2}{n_2 d_2}\right) = O\left(\frac{n_1 d_1^2}{nd_2}+ d_2\right).
\end{equation}
Since $T = O(d_1^2) = o(\norm{\vec{d}}_1)$ by~\ref{cond: stronger d_1 bound}, $k d_2 = o(\norm{\vec{d}}_1)$, and $A'=\Omega(1)$, we have
\begin{align*}
    Z &= \left(\frac{T}{2\norm{\vec{d}}_1-k d_2} - A'\cdot\frac{k d_2}{\norm{\vec{d}}_1-k d_2}\right)\left(1 + \frac{2\norm{\vec{d}}_1 - k d_2}{\norm{\vec{d}}_1 - k d_2}\cdot A' - \frac{T}{2\norm{\vec{d}}_1-k d_2}\right)
    \\
    &= \left(O\left(\frac{T}{\norm{\vec{d}}_1}\right) - O\left(\frac{d_2 A'}{\norm{\vec{d}}_1}\right)\right)\left(O\left(A'\right) - O\left(\frac{T}{\norm{\vec{d}}_1}\right)\right)
    \\
    &= O\left(\frac{T}{\norm{\vec{d}}_1} A'+\frac{T^2}{\norm{\vec{d}}_1^2}+\frac{d_2}{\norm{\vec{d}}_1} (A')^2+\frac{T d_2}{\norm{\vec{d}}_1^2} A'\right)
    \\
    &=O\left(\frac{T}{\norm{\vec{d}}_1} A'+\frac{d_2}{\norm{\vec{d}}_1} (A')^2\right),
\end{align*}
where the last equality follows from the fact that $T/\norm{\vec{d}}_1 = o(1) = o(A')$. By~\eqref{eqn: A' bound},
we have
\[
\frac{d_2 (A')^2}{\norm{\vec{d}}_1} = O\left(\frac{n_1^2 d_1^4}{n^2 d_2\norm{\vec{d}}_1} +\frac{d_2^3}{\norm{\vec{d}}_1}\right) = O\left(\frac{n_1^2 d_1^4}{n^2 d_2\norm{\vec{d}}_1} +\frac{d_1^4}{d_2\norm{\vec{d}}_1}\right) = O\left(\frac{d_1^4}{d_2\norm{\vec{d}}_1}\right) = o(1),
\]
using the fact that $n_1^2/n^2\leq n^{2\delta-2} = o(1)$ by condition~\ref{cond: n_1 small} and $d_1^4 = o(\norm{\vec{d}}_1)$ by condition~\ref{cond: stronger d_1 bound}. Furthermore, we have
\[
\frac{T}{\norm{\vec{d}}_1} A' = O\left(\frac{n_1 d_1^4}{n d_2\norm{\vec{d}}_1} + \frac{d_1^2 d_2}{\norm{\vec{d}}_1}\right) = O\left(\frac{d_1^4}{\norm{\vec{d}}_1}\right) = o(1),
\]
using the fact that $n_1/(nd_2)\leq n^{\delta-1} = o(1)$ by condition~\ref{cond: n_1 small} and $d_1^4 = o(\norm{\vec{d}}_1)$  by condition~\ref{cond: stronger d_1 bound}. Hence, $Z = o(1)$, as desired.
\end{proof}

\end{document}
